\section{Выводы}

В ходе выполнения лабораторной работы по курсу «Дискретный анализ» была спроектирована и реализована эффективная структура данных — \textbf{сбалансированное красно-чёрное дерево}. 

Основные результаты работы:
\begin{itemize}
    \item Реализованы базовые алгоритмы поддержания баланса дерева: левое и правое \textbf{вращения}, а также каскадная \textbf{перекраска} узлов, что гарантирует логарифмическую сложность основных операций.
    \item Освоена техника \textit{Pointer Tagging}: оптимизация потребления памяти путём хранения метаданных (цвета узла) в неиспользуемых младших битах указателя на родителя.
    \item Реализованы механизмы \textbf{бинарной сериализации} и десериализации, позволяющие атомарно сохранять и восстанавливать состояние дерева из внешних файлов.
    \item Проведён анализ качества кода с использованием инструмента \texttt{Valgrind} для контроля отсутствия утечек памяти и профилировщика \texttt{Callgrind} для оценки производительности критических участков алгоритма.
\end{itemize}